% Computer Science A --- JavaScript (language, DOM, async).

\runningheader{Computer Science A}{Web: JavaScript}{Page \thepage\ of \numpages}

\begingradingrange{csawebjs}

\uplevel{\par\textbf{Part 1: JavaScript fundamentals}\par}

\question[4]
You need a counter that increments on each button click. Which declaration is most appropriate for the counter binding if it will be reassigned?
\begin{choices}
  \choice \texttt{const count = 0;} (reassign later anyway)
  \CorrectChoice \texttt{let count = 0;}
  \choice \texttt{var count = 0;} only---\texttt{let} cannot be used for counters
  \choice \texttt{function count = 0;}
\end{choices}
\begin{solution}
  \texttt{let} is block-scoped and allows reassignment. \texttt{const} cannot be rebound; the bound value may still be mutable if it is an object. \texttt{var} is legacy-scoped.
\begin{lstlisting}[style=js]
let count = 0;
btn.onclick = () => { count += 1; };
\end{lstlisting}
\end{solution}

\question[4]
What happens when you access a \texttt{var} variable before its declaration line executes?
\begin{choices}
  \CorrectChoice Its value is \texttt{undefined} (due to hoisting of the declaration)
  \choice It causes a ReferenceError in all cases
  \choice It always returns \texttt{null}
  \choice The script stops at parse time with syntax error
\end{choices}
\begin{solution}
  \texttt{var} declarations are hoisted and initialized to \texttt{undefined} before execution reaches the assignment statement.
\end{solution}

\question[4]
What is the output of the code with \texttt{x = 0}, \texttt{y = '0'}, \texttt{z = false}, and this condition chain?
\begin{choices}
  \CorrectChoice \texttt{A}
  \choice \texttt{B}
  \choice \texttt{C}
  \choice \texttt{D}
\end{choices}
\begin{solution}
  Loose equality makes both comparisons true: \texttt{0 == '0'} and \texttt{'0' == false}. Therefore the first branch runs and logs \texttt{A}.
\end{solution}

\question[4]
What is logged by this loop?
\begin{lstlisting}[style=js]
let result = '';
for (let i = 5; i > 0; i--) {
  if (i % 2 === 0) {
    result += i;
  }
}
console.log(result);
\end{lstlisting}
\begin{choices}
  \CorrectChoice \texttt{42}
  \choice \texttt{24}
  \choice \texttt{54321}
  \choice \texttt{""}
\end{choices}
\begin{solution}
  Iteration visits 5,4,3,2,1 and appends only even values in that order, producing string \texttt{"42"}.
\end{solution}

\question[4]
What happens when you access a \texttt{let} variable before its declaration line executes?
\begin{choices}
  \CorrectChoice A ReferenceError occurs (temporal dead zone)
  \choice It evaluates to \texttt{undefined}
  \choice It evaluates to \texttt{null}
  \choice JavaScript silently creates a global variable
\end{choices}
\begin{solution}
  \texttt{let} and \texttt{const} bindings exist in a temporal dead zone until their declaration is evaluated.
\end{solution}

\question[4]
Which statement about JavaScript arrays is \textbf{false}?
\begin{choices}
  \CorrectChoice Arrays are fixed in size
  \choice Arrays can hold mixed value types
  \choice Arrays have a \texttt{length} property
  \choice Arrays support methods like \texttt{push} and \texttt{pop}
\end{choices}
\begin{solution}
  JavaScript arrays are dynamic and can grow or shrink.
\end{solution}

\question[4]
Which is valid JSON syntax?
\begin{choices}
  \CorrectChoice \texttt{\{"name":"John Doe","age":30,"city":"New York"\}}
  \choice \texttt{\{name:"John Doe", age:30, city:"New York"\}}
  \choice \texttt{\{name='John Doe'; age=30; city='New York'\}}
  \choice \texttt{[name: "John Doe", age: 30]}
\end{choices}
\begin{solution}
  Standard JSON requires double-quoted property names and string values.
\end{solution}

\question[4]
Which is a correct JavaScript object literal declaration?
\begin{choices}
  \CorrectChoice \texttt{let obj = \{ name: "John", age: 30 \};}
  \choice \texttt{let obj = [ name: "John", age: 30 ];}
  \choice \texttt{let obj = ( name: "John", age: 30 );}
  \choice \texttt{let obj = \{ "name" = "John", "age" = 30 \};}
\end{choices}
\begin{solution}
  Object literals use braces with key-value pairs separated by colons.
\end{solution}

\question[4]
Which of the following is \textbf{not} a falsy value in JavaScript?
\begin{choices}
  \CorrectChoice \texttt{[]}
  \choice \texttt{0}
  \choice \texttt{""}
  \choice \texttt{null}
\end{choices}
\begin{solution}
  Empty arrays are truthy, even though they may feel "empty" conceptually.
\end{solution}

\question[4]
What is the result of this code?
\begin{lstlisting}[style=js]
function addOne(value) {
  value + 1;
}
let x = 5;
addOne(x);
console.log(x);
\end{lstlisting}
\begin{choices}
  \CorrectChoice \texttt{5}
  \choice \texttt{6}
  \choice \texttt{undefined}
  \choice ReferenceError
\end{choices}
\begin{solution}
  Numbers are passed by value and the function does not return/assign the computed expression, so \texttt{x} stays 5.
\end{solution}

\question[4]
What is logged after calling \texttt{changeUsername(user)}?
\begin{lstlisting}[style=js]
function changeUsername(user) {
  user.username = 'new-username';
}
let user = { username: 'first-username' };
changeUsername(user);
console.log(user.username);
\end{lstlisting}
\begin{choices}
  \CorrectChoice \texttt{new-username}
  \choice \texttt{first-username}
  \choice \texttt{undefined}
  \choice TypeError
\end{choices}
\begin{solution}
  The object reference is shared; mutating its property inside the function updates the same object observed afterward.
\end{solution}

\question[4]
What is the result of this \texttt{switch} code when \texttt{day = 3}?
\begin{lstlisting}[style=js]
let day = 3;
let message = "";
switch (day) {
  case 1:
    message = "Monday";
  case 2:
    message += " Tuesday";
  case 3:
    message += " Wednesday";
    break;
  case 4:
    message += " Thursday";
  case 5:
    message += " Friday";
  default:
    message += " Weekend";
}
console.log(message);
\end{lstlisting}
\begin{choices}
  \CorrectChoice \texttt{Wednesday}
  \choice \texttt{Monday Tuesday Wednesday}
  \choice \texttt{Friday Weekend}
  \choice \texttt{Weekend}
\end{choices}
\begin{solution}
  Execution jumps to case 3, appends \texttt{" Wednesday"}, then breaks. No later cases run.
\end{solution}

\question[4]
What is the result of this \texttt{switch} code when \texttt{day = 5}?
\begin{lstlisting}[style=js]
let day = 5;
let message = "";
switch (day) {
  case 1:
    message = "Monday";
  case 2:
    message += " Tuesday";
  case 3:
    message += " Wednesday";
    break;
  case 4:
    message += " Thursday";
  case 5:
    message += " Friday";
  default:
    message += " Weekend";
}
console.log(message);
\end{lstlisting}
\begin{choices}
  \choice \texttt{Wednesday}
  \CorrectChoice \texttt{Friday Weekend}
  \choice \texttt{Monday Tuesday Wednesday}
  \choice \texttt{Weekend Friday}
\end{choices}
\begin{solution}
  Case 5 runs, then falls through to \texttt{default} (no \texttt{break}), so both fragments are appended.
\end{solution}

\question[4]
Which is \textbf{not} a guaranteed benefit of template literals in JavaScript?
\begin{choices}
  \choice Easier interpolation with \texttt{\$\{...\}}
  \choice Multi-line string support
  \CorrectChoice Performance improvement for every string concatenation case
  \choice Cleaner formatting for embedded expressions
\end{choices}
\begin{solution}
  Template literals improve readability and ergonomics, but they do not universally guarantee faster performance.
\end{solution}

\question[4]
What is the primary purpose of \texttt{async}/\texttt{await}?
\begin{choices}
  \CorrectChoice To handle asynchronous operations in a readable, sequential-looking style
  \choice To force all asynchronous operations to become synchronous at runtime
  \choice To remove the need for Promises internally
  \choice To speed up CPU-bound loops automatically
\end{choices}
\begin{solution}
  \texttt{await} pauses within an \texttt{async} function until a Promise settles, which simplifies async control flow.
\end{solution}

\question[4]
What effect does this CSS have?
\begin{lstlisting}[style=cssinline]
.box {
  width: 100px;
  height: 100px;
  padding: 20px;
  border: 5px solid black;
  box-sizing: border-box;
}
\end{lstlisting}
\begin{choices}
  \CorrectChoice The element's total rendered box remains 100px by 100px (padding and border included)
  \choice Content area stays 100px by 100px, making total size 150px by 150px
  \choice Border and padding are ignored because of \texttt{box-sizing}
  \choice The rule is invalid since width/height cannot mix with padding
\end{choices}
\begin{solution}
  With \texttt{box-sizing: border-box}, declared width/height include content, padding, and border.
\end{solution}

\question[4]
What is the output of \texttt{console.log(typeof typeof 1);}?
\begin{choices}
  \choice \texttt{"number"}
  \CorrectChoice \texttt{"string"}
  \choice \texttt{"undefined"}
  \choice \texttt{"object"}
\end{choices}
\begin{solution}
  \texttt{typeof 1} is the string \texttt{"number"}, and applying \texttt{typeof} again yields \texttt{"string"}.
\end{solution}

\question[4]
Given:
\begin{lstlisting}[style=htmlinline]
<p class="text" id="special">Hello World</p>
\end{lstlisting}
\begin{lstlisting}[style=cssinline]
p { color: blue; }
.text { color: green; }
#special { color: red; }
\end{lstlisting}
What color is the text?
\begin{choices}
  \choice Blue
  \choice Green
  \CorrectChoice Red
  \choice Browser default black
\end{choices}
\begin{solution}
  The ID selector \texttt{\#special} has higher specificity than class and element selectors, so red wins.
\end{solution}

\question[4]
What is logged here?
\begin{lstlisting}[style=js]
let arr = [1, 2, 3];
let [a, b, c, d] = arr;
console.log(d);
\end{lstlisting}
\begin{choices}
  \CorrectChoice \texttt{undefined}
  \choice \texttt{3}
  \choice \texttt{null}
  \choice ReferenceError
\end{choices}
\begin{solution}
  Array destructuring assigns missing positions as \texttt{undefined}.
\end{solution}

\question[4]
Given this HTML, what displays on the page?
\begin{lstlisting}[style=htmlinline]
<div>
  <p>First</p>
  <p style="display: none;">Second</p>
  <p style="visibility: hidden;">Third</p>
  <p>Fourth</p>
</div>
\end{lstlisting}
\begin{choices}
  \CorrectChoice First, then blank space where Third would be, then Fourth
  \choice First, Second, Third, Fourth
  \choice Only Fourth
  \choice First and Third only
\end{choices}
\begin{solution}
  \texttt{display: none} removes Second from layout. \texttt{visibility: hidden} hides Third but keeps its layout space.
\end{solution}

\question[4]
What is the output of:
\begin{lstlisting}[style=js]
const arr = [1, 2, 3, 4, 5];
console.log(arr.slice(1, 3));
\end{lstlisting}
\begin{choices}
  \CorrectChoice \texttt{[2, 3]}
  \choice \texttt{[1, 2, 3]}
  \choice \texttt{[2, 3, 4]}
  \choice \texttt{[3, 4]}
\end{choices}
\begin{solution}
  \texttt{slice(start, end)} includes start and excludes end.
\end{solution}

\question[4]
What are the three outputs here?
\begin{lstlisting}[style=js]
let a = 3;
let b = new Number(3);
let c = 3;
console.log(a == b);
console.log(a === b);
console.log(b === c);
\end{lstlisting}
\begin{choices}
  \CorrectChoice \texttt{true, false, false}
  \choice \texttt{true, true, true}
  \choice \texttt{false, false, true}
  \choice \texttt{false, true, false}
\end{choices}
\begin{solution}
  Loose equality coerces wrapper objects; strict equality compares type/reference, so Number object and primitive are not strictly equal.
\end{solution}

\question[4]
In the If-statement lab style, an integer \texttt{n} is \textbf{even} when which expression is true?
\begin{choices}
  \choice \texttt{n / 2 === 0}
  \CorrectChoice \texttt{n \% 2 === 0}
  \choice \texttt{n \% 2 === 1}
  \choice \texttt{Math.floor(n) === n / 2}
\end{choices}
\begin{solution}
  The \textbf{remainder} operator \texttt{\%} yields zero for even \texttt{n} under \texttt{n \% 2 === 0}. Odd numbers leave remainder \texttt{1}.
\begin{lstlisting}[style=js]
if (n % 2 === 0) {
  console.log("even");
}
\end{lstlisting}
\end{solution}

\question[4]
In \texttt{if (cond1) \{ A \} else if (cond2) \{ B \} else \{ C \}}, block \texttt{C} runs when:
\begin{choices}
  \choice \texttt{cond1} and \texttt{cond2} are both truthy
  \CorrectChoice \texttt{cond1} is falsy and \texttt{cond2} is falsy
  \choice \texttt{cond1} is truthy regardless of \texttt{cond2}
  \choice Only when \texttt{cond2} throws an error
\end{choices}
\begin{solution}
  The chain tests \texttt{cond1} then, if it failed, \texttt{cond2}. The \texttt{else} branch runs only when every tested condition above was falsy (If-statement lab structure).
\begin{lstlisting}[style=js]
if (cond1) {
  A();
} else if (cond2) {
  B();
} else {
  C();
}
\end{lstlisting}
\end{solution}

\newpage

\question[4]
In \texttt{for (init; condition; final)}, the \texttt{final} clause (for example \texttt{i++}) runs:
\begin{choices}
  \choice Once before any iteration
  \choice After the body on each iteration
  \choice Only if the loop body throws
  \CorrectChoice After each iteration body completes, before the next condition check
\end{choices}
\begin{solution}
  A C-style \texttt{for} loop runs \textbf{init} once, then each cycle: test \textbf{condition}, run body, run \textbf{final}, repeat. That matches the lab's three-part loop table.
\begin{lstlisting}[style=js]
for (let i = 0; i < 3; i++) {
  console.log(i);
}
\end{lstlisting}
\end{solution}

\question[4]
Which value is \textbf{falsy} in JavaScript when used as an \texttt{if} condition?
\begin{choices}
  \choice \texttt{[]} (empty array)
  \choice \texttt{\{\}} (empty object)
  \CorrectChoice \texttt{""} (empty string)
  \choice \texttt{"0"}
\end{choices}
\begin{solution}
  The short falsy list includes \texttt{false}, \texttt{0}, \texttt{-0}, \texttt{0n}, \texttt{""}, \texttt{null}, \texttt{undefined}, \texttt{NaN}. Empty arrays and objects are \textbf{truthy}---a common beginner surprise from the Truthyfalsy lab.
\begin{lstlisting}[style=js]
if ("") { /* skipped */ }
if ([]) { /* runs: [] is truthy */ }
\end{lstlisting}
\end{solution}

\question[4]
What is the result of \texttt{5 === "5"} in JavaScript?
\begin{choices}
  \CorrectChoice \texttt{false} (no type coercion with strict equality)
  \choice \texttt{true} because both represent five
  \choice \texttt{undefined}
  \choice A syntax error
\end{choices}
\begin{solution}
  \textbf{Strict equality} \texttt{===} compares \textbf{without} coercion, so number and string differ. \texttt{5 == "5"} is \texttt{true} with loose equality (coercion). Modern style prefers \texttt{===} / \texttt{!==}.
\begin{lstlisting}[style=js]
console.log(5 === "5"); // false
\end{lstlisting}
\end{solution}


\newpage
\question[4]
What is the value of \texttt{5 == "5"} (loose equality) in JavaScript?
\begin{choices}
  \CorrectChoice \texttt{true} (operands may be coerced before comparison)
  \choice \texttt{false}
  \choice \texttt{"55"}
  \choice It throws \texttt{TypeError}
\end{choices}
\begin{solution}
  Loose \texttt{==} applies the language's coercion rules; numeric string \texttt{"5"} converts to number \texttt{5}. The Typecoercion lab contrasts this with predictable \texttt{===}.
\begin{lstlisting}[style=js]
console.log(5 == "5"); // true
\end{lstlisting}
\end{solution}

\question[4]
What does \texttt{Number("42")} evaluate to?
\begin{choices}
  \CorrectChoice The number \texttt{42}
  \choice The string \texttt{"42"}
  \choice \texttt{NaN}
  \choice \texttt{true}
\end{choices}
\begin{solution}
  \texttt{Number(...)} performs \textbf{explicit} numeric conversion. Invalid strings yield \texttt{NaN}; the lab contrasted this with implicit coercion in \texttt{==} comparisons.
\begin{lstlisting}[style=js]
const x = Number("42"); // 42
\end{lstlisting}
\end{solution}

\question[4]
For an array \texttt{arr}, \texttt{arr.push("x")} primarily:
\begin{choices}
  \choice Removes the last element
  \CorrectChoice Adds an element at the end and returns the new \texttt{length}
  \choice Sorts the array in place
  \choice Returns the pushed value only, without changing \texttt{arr}
\end{choices}
\begin{solution}
  \texttt{push} appends; \texttt{pop} removes from the end. \texttt{length} tracks element count (for dense arrays, last index is \texttt{length - 1}).
\begin{lstlisting}[style=js]
const arr = [1, 2];
arr.push("x"); // length is now 3
\end{lstlisting}
\end{solution}

\question[4]
Given \texttt{let user = \{ name: "Ada", score: 10 \};}, which expression reads the score using \textbf{bracket} notation with a string key?
\begin{choices}
  \choice \texttt{user.score.}
  \CorrectChoice \texttt{user["score"]}
  \choice \texttt{user::score}
  \choice \texttt{user(score)}
\end{choices}
\begin{solution}
  \textbf{Dot} syntax (\texttt{user.score}) is idiomatic for valid identifiers. \textbf{Brackets} support dynamic keys and non-identifier keys (\texttt{user["score"]}).
\begin{lstlisting}[style=js]
const user = { name: "Ada", score: 10 };
const s = user["score"];
\end{lstlisting}
\end{solution}

\question[4]
In the Button-click lab style, which API retrieves a single element by its \texttt{id} attribute?
\begin{choices}
  \choice \texttt{document.querySelector("id")}
  \CorrectChoice \texttt{document.getElementById("...")}
  \choice \texttt{document.getElementsByTagName("id")}
  \choice \texttt{window.findId("...")}
\end{choices}
\begin{solution}
  \texttt{getElementById} returns one \texttt{Element} or \texttt{null}. The lab pairs this with \texttt{onclick} assignment or later \texttt{addEventListener} patterns.
\begin{lstlisting}[style=js]
const btn = document.getElementById("go");
\end{lstlisting}
\end{solution}

\question[4]
To replace all text inside a DOM element \texttt{el} with the string \texttt{"Done"}, you should assign to:
\begin{choices}
  \choice \texttt{el.innerHTML = "Done"} only (always safest)
  \CorrectChoice \texttt{el.textContent = "Done"} (plain text, avoids parsing HTML)
  \choice \texttt{el.value = "Done"} on any element
  \choice \texttt{el.node = "Done"}
\end{choices}
\begin{solution}
  \texttt{textContent} sets character data without interpreting markup. \texttt{innerHTML} parses HTML (XSS risk with untrusted strings). \texttt{.value} applies to form controls.
\begin{lstlisting}[style=js]
el.textContent = "Done";
\end{lstlisting}
\end{solution}

\question[4]
To read what the user currently typed in a text \texttt{<input>} element referenced as \texttt{inputEl}, use:
\begin{choices}
  \choice \texttt{inputEl.textContent}
  \CorrectChoice \texttt{inputEl.value}
  \choice \texttt{inputEl.getAttribute("value")} always matches live typing
  \choice \texttt{inputEl.innerText} only
\end{choices}
\begin{solution}
  The live string for most inputs is \texttt{.value}. The HTML \texttt{value} attribute is only the \emph{initial} default unless you sync it; the Input lab stresses reading \texttt{.value} for current text.
\begin{lstlisting}[style=js]
const text = inputEl.value;
\end{lstlisting}
\end{solution}

\question[4]
What is the value of \texttt{typeof 42} in JavaScript?
\begin{choices}
  \choice \texttt{"integer"}
  \choice \texttt{"float"}
  \CorrectChoice \texttt{"number"}
  \choice \texttt{"decimal"}
\end{choices}
\begin{solution}
  JavaScript has one numeric type in \texttt{typeof} results: \texttt{"number"} (both \texttt{3} and \texttt{3.14}).
\begin{lstlisting}[style=js]
typeof 42;        // "number"
typeof "42";       // "string"
\end{lstlisting}
\end{solution}

\question[4]
What does \texttt{"5" + 1} evaluate to?
\begin{choices}
  \choice \texttt{6} (numeric addition)
  \choice \texttt{NaN}
  \CorrectChoice \texttt{"51"} (string concatenation after coercion)
  \choice The string \texttt{"6"}
\end{choices}
\begin{solution}
  With \texttt{+}, if either operand is a string, the other is typically coerced to string and concatenated---not added numerically.
\begin{lstlisting}[style=js]
console.log("5" + 1);   // "51"
console.log("5" - 1);   // 4  (subtraction uses numeric coercion)
\end{lstlisting}
\end{solution}

\question[4]
Which expression uses a \textbf{template literal} to embed \texttt{name} into a greeting?
\begin{choices}
  \choice \texttt{'Hello, ' + name + '!'}
  \CorrectChoice \texttt{`Hello, \$\{name\}!`}
  \choice \texttt{"Hello, \#name!"}
  \choice \texttt{concat("Hello, ", name)}
\end{choices}
\begin{solution}
  Backtick strings allow \texttt{\$\{\ldots\}} interpolations; they are the ES6 template literal syntax.
\begin{lstlisting}[style=js]
const name = "Ada";
const msg = `Hello, ${name}!`;
\end{lstlisting}
\end{solution}

\question[4]
\texttt{JSON.parse('\{"score":10\}')} returns:
\begin{choices}
  \choice A string equal to \texttt{"score"}
  \CorrectChoice A JavaScript object (e.g.\ you can read \texttt{.score})
  \choice \texttt{undefined}
  \choice A syntax error because single quotes are illegal
\end{choices}
\begin{solution}
  \texttt{JSON.parse} deserializes a JSON text into JS values (objects, arrays, strings, numbers, booleans, null). \texttt{JSON.stringify} does the reverse.
\begin{lstlisting}[style=js]
const obj = JSON.parse('{"score":10,"ok":true}');
console.log(obj.score);
\end{lstlisting}
\end{solution}

\question[4]
Compared to \texttt{splice}, \texttt{Array.prototype.slice} is different because \texttt{slice}:
\begin{choices}
  \choice Always empties the original array
  \choice Only works on strings, not arrays
  \CorrectChoice Returns a shallow copy of a section and does \textbf{not} mutate the original array
  \choice Sorts the array in descending order
\end{choices}
\begin{solution}
  \texttt{slice(start, end)} copies a range; \texttt{splice} changes the array in place (insert/remove). Recognizing mutating vs non-mutating methods prevents subtle bugs.
\begin{lstlisting}[style=js]
const nums = [10, 20, 30];
const part = nums.slice(0, 2); // [10, 20], nums unchanged
\end{lstlisting}
\end{solution}

\question[4]
The operator \texttt{!==} is best described as:
\begin{choices}
  \choice Loose inequality (allows type coercion)
  \CorrectChoice Strict inequality (no type coercion, opposite of \texttt{===})
  \choice Bitwise NOT combined with assignment
  \choice Legal only inside \texttt{for} loops
\end{choices}
\begin{solution}
  Pair \texttt{!==} with \texttt{===} for predictable comparisons; \texttt{!=} is the loose form paired with \texttt{==}.
\begin{lstlisting}[style=js]
5 !== "5"; // true
5 !=  "5"; // false
\end{lstlisting}
\end{solution}

\question[4]
For a non-empty array \texttt{arr}, the index of the \textbf{last} element is:
\begin{choices}
  \choice \texttt{arr.length}
  \CorrectChoice \texttt{arr.length - 1}
  \choice \texttt{arr.length + 1}
  \choice Always \texttt{0}
\end{choices}
\begin{solution}
  Arrays are zero-indexed: first at \texttt{0}, last at \texttt{length - 1}. Off-by-one errors at \texttt{length} are a common debugging target.
\begin{lstlisting}[style=js]
const arr = ["a", "b", "c"];
const last = arr[arr.length - 1];
\end{lstlisting}
\end{solution}

\question[4]
What are the results of \texttt{3 == '3'} and \texttt{3 === '3'} in JavaScript?
\begin{choices}
  \CorrectChoice \texttt{true} then \texttt{false}---\texttt{==} may coerce types; \texttt{===} does not
  \choice \texttt{false} then \texttt{true}
  \choice Both are \texttt{false}
  \choice Both are \texttt{true}
\end{choices}
\begin{solution}
  Loose equality coerces the string \texttt{"3"} to the number \texttt{3}. Strict equality compares type and value, so number and string differ.
\begin{lstlisting}[style=js]
console.log(3 == "3");  // true
console.log(3 === "3"); // false
\end{lstlisting}
\end{solution}

\question[4]
Which keywords can declare a variable binding in modern JavaScript?
\begin{choices}
  \CorrectChoice \texttt{let}, \texttt{var}, and \texttt{const}
  \choice \texttt{int} and \texttt{async}
  \choice Only \texttt{var} (\texttt{let} and \texttt{const} are invalid)
  \choice \texttt{for} and \texttt{function} only
\end{choices}
\begin{solution}
  \texttt{let} and \texttt{const} are block-scoped; \texttt{var} is function-scoped (legacy). \texttt{const} forbids rebinding but object contents may still mutate.
\end{solution}

\question[4]
JavaScript (ES6) has separate built-in \texttt{int} and \texttt{double} numeric types.
\begin{choices}
  \choice True
  \CorrectChoice False---there is one \texttt{number} type (IEEE-754); \texttt{typeof} reports \texttt{"number"}
  \choice True only in strict mode
  \choice True only in browsers, not in Node
\end{choices}
\begin{solution}
  All numeric literals share the same type; distinguish integers with \texttt{Number.isInteger} when needed.
\end{solution}

\question[4]
Which line correctly creates an empty \textbf{object literal}?
\begin{choices}
  \CorrectChoice \texttt{let myObject = \{\};}
  \choice \texttt{let Object myObject = new Object();}
  \choice \texttt{let myObject = ();}
  \choice JavaScript has no object type
\end{choices}
\begin{solution}
  Curly braces with optional properties define object literals; \texttt{new Object()} works but the literal form is idiomatic.
\begin{lstlisting}[style=js]
let myObject = {};
myObject.score = 10;
\end{lstlisting}
\end{solution}

\question[4]
How do you read the \textbf{first} element of an array \texttt{arr}?
\begin{choices}
  \choice \texttt{arr(0)}
  \CorrectChoice \texttt{arr[0]}
  \choice \texttt{arr[1]}
  \choice \texttt{arr\{0\}}
\end{choices}
\begin{solution}
  Arrays are zero-indexed: first at \texttt{0}, last at \texttt{length - 1}.
\end{solution}

\question[4]
You need a template literal that interpolates \texttt{product.quantity}, \texttt{product.name}, and \texttt{product.price}, and prints a literal dollar sign immediately before the numeric price. Which pattern is correct?
\begin{choices}
  \choice Backticks with \texttt{\$product.quantity} (no braces)
  \CorrectChoice Backticks with \texttt{\$\{product.quantity\}}, \texttt{\$\{product.name\}}, and \texttt{\$\$\{product.price\}} (first \texttt{\$} is literal, \texttt{\$\{...\}} interpolates)
  \choice Single quotes with \texttt{\$\{\{product.quantity\}\}}
  \choice Single quotes with \texttt{\$\{product.quantity\}}
\end{choices}
\begin{solution}
  Inside template literals, \texttt{\$\{expr\}} evaluates an expression. To output a literal \texttt{\$}, escape by writing \texttt{\$\$} before \texttt{\$\{product.price\}}.
\begin{lstlisting}[style=js]
const s = `${product.quantity} units of ${product.name} at $${product.price}`;
\end{lstlisting}
\end{solution}

\question[4]
What happens when this code runs?
\begin{lstlisting}[style=js]
const a = 0;
for (let i = 5; i < 10; i++) {
  a++;
}
\end{lstlisting}
\begin{choices}
  \choice \texttt{a} is \texttt{5}
  \choice \texttt{a} is \texttt{4}
  \choice \texttt{a} is \texttt{undefined}
  \CorrectChoice A \textbf{TypeError}: \texttt{const} bindings cannot be reassigned (\texttt{a++} mutates the binding)
\end{choices}
\begin{solution}
  \texttt{const} fixes the binding to the value \texttt{0}; increment operators require a mutable binding---use \texttt{let} for counters.
\end{solution}

\newpage
\uplevel{\par\textbf{Part 2: DOM manipulation \& events}\par}

\question[4]
To change which image file an \texttt{<img>} displays, you can assign to \texttt{imgEl.src} or call:
\begin{choices}
  \choice \texttt{imgEl.href = "pic.png"}
  \CorrectChoice \texttt{imgEl.setAttribute("src", "pic.png")}
  \choice \texttt{imgEl.alt = "pic.png"}
  \choice \texttt{imgEl.createElement("src")}
\end{choices}
\begin{solution}
  The Attributes lab uses \texttt{.src} or \texttt{setAttribute("src", ...)} so the browser loads a new resource. \texttt{alt} is alternative text, not the image URL.
\begin{lstlisting}[style=js]
imgEl.setAttribute("src", "pic.png");
/* or: imgEl.src = "pic.png"; */
\end{lstlisting}
\end{solution}

\question[4]
To add a new \texttt{<li>} to an existing \texttt{<ol>} in the Appending-elements lab pattern, you typically:
\begin{choices}
  \choice Call \texttt{ol.innerHTML += "<li>"} only (no creation step)
  \CorrectChoice \texttt{document.createElement("li")} then \texttt{parent.appendChild(newLi)}
  \choice \texttt{ol.push(li)}
  \choice \texttt{new Element("li")}
\end{choices}
\begin{solution}
  \texttt{createElement} builds a node; \texttt{appendChild} attaches it under a parent. String \texttt{innerHTML} reparses markup and can drop event listeners---\texttt{createElement} is the structured approach.
\begin{lstlisting}[style=js]
const li = document.createElement("li");
li.textContent = "New item";
ol.appendChild(li);
\end{lstlisting}
\end{solution}

\question[4]
Compared to assigning \texttt{element.onclick = fn}, \texttt{element.addEventListener("click", fn)} is especially useful because it:
\begin{choices}
  \choice Runs the handler before the event exists
  \CorrectChoice Can register \textbf{multiple} listeners for the same event without overwriting earlier ones
  \choice Only works for keyboard events
  \choice Prevents event bubbling always
\end{choices}
\begin{solution}
  \texttt{onclick} property assignment replaces the previous handler. \texttt{addEventListener} stacks listeners and offers capture/bubbling control (\texttt{Eventlistener} lab).
\begin{lstlisting}[style=js]
el.addEventListener("click", fn1);
el.addEventListener("click", fn2);
\end{lstlisting}
\end{solution}

\newpage

\question[4]
Inside a click handler, \texttt{event.target} typically refers to:
\begin{choices}
  \choice The \texttt{window} object always
  \CorrectChoice The innermost element that was the origin of the event
  \choice The document root always
  \choice The function that handled the event
\end{choices}
\begin{solution}
  The \textbf{Event} object reports \texttt{type}, \texttt{target} (where the event started), and related fields. The Eventobject lab used \texttt{target} after a click.
\begin{lstlisting}[style=js]
el.addEventListener("click", (event) => {
  console.log(event.target);
});
\end{lstlisting}
\end{solution}

\question[4]
\texttt{document.querySelector("\#go")} differs from \texttt{document.getElementById("go")} mainly in that \texttt{querySelector}:
\begin{choices}
  \choice Can only match tags, never IDs
  \CorrectChoice Takes a full \textbf{CSS selector string} and returns the \textbf{first} matching element (or \texttt{null})
  \choice Always returns a NodeList of all matches
  \choice Throws if the ID does not exist
\end{choices}
\begin{solution}
  \texttt{querySelector} is the general CSS-selector API; prefix IDs with \texttt{\#}, classes with \texttt{.}, etc. \texttt{getElementById} is ID-only and very slightly older/simpler.
\begin{lstlisting}[style=js]
const a = document.querySelector("#go");
const b = document.getElementById("go");
\end{lstlisting}
\end{solution}

\question[4]
Calling \texttt{event.preventDefault()} inside a click handler on a link typically:
\begin{choices}
  \choice Deletes the link from the DOM
  \CorrectChoice Stops the browser's default action (e.g.\ following the \texttt{href}) while still allowing the handler to run
  \choice Removes all other listeners on that element
  \choice Is required for every \texttt{addEventListener} call
\end{choices}
\begin{solution}
  \texttt{preventDefault} blocks default behavior (navigation, form submit, context menu, etc.). It is unrelated to \texttt{stopPropagation}, which affects event bubbling/capture.
\begin{lstlisting}[style=js]
link.addEventListener("click", (event) => {
  event.preventDefault();
  console.log("No navigation");
});
\end{lstlisting}
\end{solution}

\newpage
\question[4]
To add a CSS class \texttt{active} without overwriting existing classes on \texttt{el}, the preferred DOM API is:
\begin{choices}
  \choice \texttt{el.class = "active"}
  \choice \texttt{el.className += "active"} only (no space)
  \CorrectChoice \texttt{el.classList.add("active")}
  \choice \texttt{el.setAttribute("classList", "active")}
\end{choices}
\begin{solution}
  \texttt{classList} maintains the class token set: \texttt{add}, \texttt{remove}, \texttt{toggle}, \texttt{contains}. Assigning \texttt{className} replaces the whole attribute string.
\begin{lstlisting}[style=js]
el.classList.add("active");
el.classList.remove("hidden");
\end{lstlisting}
\end{solution}

\question[4]
For \texttt{<div id="app" data-user-id="7">}, the JavaScript property that reads \texttt{"7"} is:
\begin{choices}
  \choice \texttt{el.data-user-id} (valid property syntax)
  \choice \texttt{el.getAttribute("user-id")}
  \CorrectChoice \texttt{el.dataset.userId}
  \choice \texttt{el.dataset.data-user-id}
\end{choices}
\begin{solution}
  \texttt{data-*} attributes appear on \texttt{element.dataset} in camelCase: \texttt{data-user-id} $\rightarrow$ \texttt{dataset.userId}. Values are always strings.
\begin{lstlisting}[style=htmlinline]
<div id="app" data-user-id="7" data-role="admin"></div>
\end{lstlisting}
\begin{lstlisting}[style=js]
const el = document.getElementById("app");
console.log(el.dataset.userId); // "7"
\end{lstlisting}
\end{solution}

\question[4]
To remove a child node \texttt{child} from its parent \texttt{parent}, modern code often uses:
\begin{choices}
  \choice \texttt{parent.delete(child)}
  \choice \texttt{child.erase()}
  \CorrectChoice \texttt{child.remove()} (or \texttt{parent.removeChild(child)})
  \choice \texttt{document.unlink(child)}
\end{choices}
\begin{solution}
  \texttt{Element.remove()} detaches the node from the tree. The older \texttt{removeChild} pattern still appears in legacy code but is more verbose.
\begin{lstlisting}[style=js]
child.remove();
/* or: parent.removeChild(child); */
\end{lstlisting}
\end{solution}

\newpage

\uplevel{\par\textbf{Part 3: Advanced JavaScript}\par}

\question[4]
In an ES6 class, which special method runs when you write \texttt{new MyClass(args)}?
\begin{choices}
  \choice \texttt{static init()}
  \CorrectChoice \texttt{constructor(args) \{ \ldots \}}
  \choice \texttt{function MyClass()}
  \choice \texttt{this.build()}
\end{choices}
\begin{solution}
  \texttt{constructor} initializes the instance; \texttt{this} refers to that instance inside instance methods. The Classes lab covered instantiation and calling methods like \texttt{getDescription()}.
\begin{lstlisting}[style=js]
class Widget {
  constructor(name) { this.name = name; }
}
const w = new Widget("clock");
\end{lstlisting}
\end{solution}

\question[4]
Which construct runs \texttt{recover()} only if \texttt{risky()} throws?
\begin{choices}
  \choice \texttt{if (risky()) recover();}
  \CorrectChoice \texttt{try \{ risky(); \} catch (e) \{ recover(); \}}
  \choice \texttt{throw risky() || recover();}
  \choice \texttt{risky().catch(recover)} for sync code only
\end{choices}
\begin{solution}
  \texttt{try}/\texttt{catch} handles synchronous exceptions; \texttt{throw} signals an error. The Errors lab stressed preventing uncaught exceptions from crashing the script path under test.
\begin{lstlisting}[style=js]
try { risky(); } catch (e) { recover(); }
\end{lstlisting}
\end{solution}

\question[4]
\texttt{setTimeout(callback, 1000)} schedules the \texttt{callback} to run:
\begin{choices}
  \choice Immediately in the same call stack
  \CorrectChoice After roughly 1000\,ms, via the browser's timer queue (async relative to current synchronous code)
  \choice Only if \texttt{callback} returns \texttt{true}
  \choice Before any \texttt{console.log} on the page
\end{choices}
\begin{solution}
  \textbf{Callbacks} passed to \texttt{setTimeout} run later; order relative to other macrotasks depends on the event loop (Callbackfunction lab).
\begin{lstlisting}[style=js]
setTimeout(() => console.log("later"), 1000);
\end{lstlisting}
\end{solution}

\newpage

\question[4]
Declaring a function as \texttt{async function f() \{ \ldots \}} means that calling \texttt{f()} returns:
\begin{choices}
  \choice A plain number or string always
  \choice \texttt{undefined} always
  \CorrectChoice A \textbf{Promise} (resolved with the function's returned value or rejected on throw)
  \choice A synchronous blocking lock
\end{choices}
\begin{solution}
  \texttt{async} functions wrap the return value in a Promise. \texttt{await} pauses \emph{inside} the async function until a Promise settles (Asyncawait lab).
\begin{lstlisting}[style=js]
async function f() { return 1; }
const p = f(); // a Promise
\end{lstlisting}
\end{solution}

\question[4]
The expression \texttt{await fetch(url)} (inside an \texttt{async} function) typically yields:
\begin{choices}
  \choice Parsed JSON immediately
  \CorrectChoice A \texttt{Response} object (body still unread until \texttt{json()}, \texttt{text()}, etc.)
  \choice A string URL
  \choice An \texttt{ArrayBuffer} always
\end{choices}
\begin{solution}
  \texttt{fetch} resolves to a \texttt{Response}; you then \texttt{await res.json()} or similar to parse. The Fetch lab chained request $\rightarrow$ parse $\rightarrow$ DOM update.
\begin{lstlisting}[style=js]
async function load(url) {
  const res = await fetch(url);
  const data = await res.json();
  return data;
}
\end{lstlisting}
\end{solution}

\question[4]
In \texttt{Math.max(...values)}, the \texttt{...values} syntax in the call is called:
\begin{choices}
  \choice A rest parameter
  \CorrectChoice Spread (expands an iterable into individual arguments)
  \choice A template literal
  \choice Destructuring assignment
\end{choices}
\begin{solution}
  \textbf{Spread} in a call expands an array (or iterable) into arguments. \textbf{Rest} in the parameter list (\texttt{function f(...args)}) collects extra arguments into an array (Spread/rest lab).
\begin{lstlisting}[style=js]
const values = [3, 9, 4];
Math.max(...values);
\end{lstlisting}
\end{solution}


\newpage
\question[4]
In \texttt{function sum(...nums) \{ return nums.reduce((a,b)=>a+b,0); \}}, the \texttt{...nums} in the parameter list is:
\begin{choices}
  \choice Spread syntax
  \CorrectChoice Rest syntax (collects remaining arguments into an array \texttt{nums})
  \choice An object spread
  \choice Illegal in strict mode
\end{choices}
\begin{solution}
  Rest parameters bundle trailing arguments. Spread and rest share the \texttt{...} token but appear in opposite ``directions'' (call vs.\ parameter list).
\begin{lstlisting}[style=js]
function sum(...nums) {
  return nums.reduce((a, b) => a + b, 0);
}
\end{lstlisting}
\end{solution}

\question[4]
After \texttt{const res = await fetch(url);}, a robust pattern is to check \texttt{res.ok} because:
\begin{choices}
  \choice \texttt{fetch} throws automatically on HTTP 404
  \CorrectChoice \texttt{fetch} resolves even for error status codes (4xx/5xx); \texttt{res.ok} is false when the status is not in the 200--299 range
  \choice \texttt{res.ok} is true only before \texttt{await}
  \choice \texttt{res.ok} replays the \texttt{network\_request}
\end{choices}
\begin{solution}
  Failed HTTP responses still yield a \texttt{Response}; inspect \texttt{status} or the boolean \texttt{ok} before calling \texttt{json()}.
\begin{lstlisting}[style=js]
const res = await fetch(url);
if (!res.ok) throw new Error(`HTTP ${res.status}`);
const data = await res.json();
\end{lstlisting}
\end{solution}

\question[4]
If an \texttt{async} function body executes \texttt{throw new Error("bad");}, the Promise returned by calling that function will:
\begin{choices}
  \choice Stay pending forever
  \choice Always resolve to \texttt{undefined}
  \CorrectChoice Become \textbf{rejected} with that error (unless caught inside the function)
  \choice Throw synchronously in the caller without a Promise
\end{choices}
\begin{solution}
  Exceptions inside \texttt{async} functions reject the implicit Promise. Callers should use \texttt{try}/\texttt{catch} with \texttt{await} or \texttt{.catch} on the Promise.
\begin{lstlisting}[style=js]
async function boom() { throw new Error("bad"); }
boom().catch((e) => console.log(e.message));
\end{lstlisting}
\end{solution}

\newpage

\question[4]
\texttt{Promise.all([p1, p2])} resolves when:
\begin{choices}
  \choice The \textbf{first} promise settles (either fulfilled or rejected)
  \CorrectChoice \textbf{Every} promise in the array has fulfilled; if any rejects, the whole \texttt{Promise.all} rejects
  \choice All promises are still pending
  \choice Only when called inside a \texttt{for} loop
\end{choices}
\begin{solution}
  \texttt{Promise.all} is useful for parallel async work; combine with \texttt{Promise.allSettled} when you need all outcomes regardless of failures (ES2020+).
\begin{lstlisting}[style=js]
const [a, b] = await Promise.all([fetch(url1), fetch(url2)]);
\end{lstlisting}
\end{solution}

\question[4]
The method \texttt{res.json()} on a \texttt{fetch} \texttt{Response} is \textbf{asynchronous} because:
\begin{choices}
  \choice It always returns \texttt{undefined}
  \choice It requires two URLs
  \CorrectChoice Reading/parsing the body is streamed and completes later---it returns a \textbf{Promise}
  \choice JSON is only valid in Node.js, not browsers
\end{choices}
\begin{solution}
  Body consumption methods (\texttt{json}, \texttt{text}, \texttt{blob}) are async; you must \texttt{await} (or use \texttt{.then}) even after you already have \texttt{res}.
\begin{lstlisting}[style=js]
const res = await fetch(url);
const data = await res.json();
\end{lstlisting}
\end{solution}

\question[4]
Given:
\begin{lstlisting}[style=js]
let promise = new Promise((resolve, reject) => {
  reject("Promise rejected!");
});
promise.catch((error) => console.log(error));
promise.catch((error) => console.log(error));
\end{lstlisting}
What is printed?
\begin{choices}
  \choice Nothing
  \CorrectChoice \texttt{Promise rejected!} appears \textbf{twice} (each \texttt{.catch} on the same promise can handle the rejection)
  \choice \texttt{undefined} twice
  \choice Only an unhandled rejection warning
\end{choices}
\begin{solution}
  Registering multiple \texttt{.catch} handlers on one promise means each handler may run for the same rejection (they are separate registrations on the same settlement).
\end{solution}

\question[4]
In JavaScript, function values are objects and inherit (via the prototype chain) from \texttt{Object}.
\begin{choices}
  \CorrectChoice True
  \choice False
  \choice True only for arrow functions
  \choice True only for class methods
\end{choices}
\begin{solution}
  Functions are callable objects; their prototypes ultimately link like other objects, which is why you can attach properties to a function.
\end{solution}

\question[4]
What is \textbf{event capturing} in the DOM event model?
\begin{choices}
  \CorrectChoice The event is dispatched from the root down through ancestors toward the target (\textbf{capture} phase) before the target and bubble phases
  \choice The event propagates only from the target upward (bubbling only)
  \choice The browser logs every event to the console automatically
  \choice Events jump between sibling nodes at random
\end{choices}
\begin{solution}
  With \texttt{addEventListener(..., true)}, handlers run during capture; the default bubble phase goes target $\rightarrow$ ancestors. \texttt{stopPropagation} affects both.
\end{solution}

\question[4]
Which kind of JavaScript value cannot be represented in standard JSON text?
\begin{choices}
  \CorrectChoice A function value
  \choice A string
  \choice A plain object
  \choice An array
\end{choices}
\begin{solution}
  JSON supports strings, numbers, booleans, \texttt{null}, arrays, and objects. Functions, \texttt{undefined}, and symbols are omitted or require custom serializers.
\end{solution}

\question[4]
Which API serializes a JavaScript value into a JSON \textbf{string}?
\begin{choices}
  \choice \texttt{JSON.toString()}
  \CorrectChoice \texttt{JSON.stringify()}
  \choice \texttt{String.JSON()}
  \choice \texttt{json.toString()}
\end{choices}
\begin{solution}
  Pair \texttt{JSON.stringify} with \texttt{JSON.parse} for round-trips; handle replacers/revivers when dates or special values matter.
\begin{lstlisting}[style=js]
const s = JSON.stringify({ score: 10 });
const obj = JSON.parse(s);
\end{lstlisting}
\end{solution}

\question[4]
The browser \texttt{fetch} API is primarily used to \ldots
\begin{choices}
  \choice Read object fields without referencing the object
  \choice Force any Promise to resolve synchronously
  \CorrectChoice Perform HTTP requests and obtain \texttt{Response} objects (often followed by \texttt{json()} or \texttt{text()})
  \choice Search lexical scope for a variable name
\end{choices}
\begin{solution}
  \texttt{fetch} is the modern network entry point; it returns a Promise of a \texttt{Response} and does not throw on HTTP error status codes by default.
\end{solution}

\endgradingrange{csawebjs}
